\documentclass[10pt,a4paper]{amsart}
\usepackage[margin=19mm]{geometry}
\usepackage{amsmath,amssymb,mathtools}
\usepackage[scr=rsfs,cal=euler]{mathalfa}
\usepackage{xfrac}
\usepackage[unicode,hidelinks,bookmarks=false]{hyperref}
\newcommand{\ie}{i.e.\ }
\newcommand{\cf}{cf.\ }
\newcommand{\resp}{respectively\ }
\newcommand{\Subsection}{\subsection}
\providecommand{\nobreakdash}{\nobreak\hbox{-}}
\setlength{\emergencystretch}{2em}
\multlinegap0pt
\usepackage{mathtools}
\usepackage{booktabs}
\usepackage{amssymb}
\usepackage[shortlabels]{enumitem}
\usepackage{xcolor}
\usepackage{cancel}
\usepackage{tikz}
\newtheorem{lemma}{Lemma}
\DeclareMathOperator{\Ker}{Ker}
\DeclareMathOperator{\Ran}{Ran}
\title{Frames generated by the functional calculus and function frames of a normal operator}
\author{Nizar El Idrissi}
\date{Source: arXiv:2107.10290v5 (CC BY 4.0)}
\begin{document}
\maketitle

\noindent\emph{Source and licence.} Frames generated by the functional calculus and function frames of a normal operator. Version arXiv:2107.10290v5 (CC BY 4.0), distributed by arXiv under the Creative Commons Attribution 4.0 International licence, http://creativecommons.org/licenses/by/4.0/, as recorded in the arXiv licence metadata for this exact version and shown on the abstract page https://arxiv.org/abs/2107.10290v5. Full licence notice: \texttt{licenses/arxiv-2107.10290-CC-BY-4.0.txt}.\par\smallskip

\noindent\textit{Editorial scope.} Counted unit: the article's lemma Lemma(T-L)SurjApproxSpecAdjoint (source0.tex lines 440--444). The article's complete original proof of it is delivered with it (source0.tex lines 446--499, one \texttt{proof} environment of 54 lines). No mathematics is written, repaired or re-derived: the statement and the proof are the article's own lines, in the article's own notation, and no retained line is substituted or re-flowed.  No further source-local context is needed: every cross-reference of every retained line resolves inside the two delivered spans.  Nothing was omitted from any retained span, so the package carries no omission record.  No retained line contains a citation, so the wrapper carries no bibliography and no bibliography entry.  No retained line contains a figure environment, an image-inclusion command or drawing code, so no image is delivered and the package declares no asset dependency.  Editorial formatting only, in addition: an AMS \texttt{amsart} preamble that restates the article's own declarations and macros used by the retained lines, namely the environments \texttt{lemma} on separate counters, together with the standard packages those lines need.  The retained environments print in this document as Lemma 1 in order of appearance, whereas the article prints the same items with the numbering of its own full document.  The complete native source of the article as distributed in the arXiv e-print for this exact version is bundled unchanged as \texttt{sources/114087/source0.tex}.
\begin{lemma}[Surjectivity and approximate point spectrum]
\label{Lemma(T-L)SurjApproxSpecAdjoint}
Let $K$ be a Hilbert space, $T \in \mathcal{B}(K)$, and $\lambda \in \mathbb{C}$. Then:
\[\overline{\lambda} \notin \sigma_{ap}(T^*) \Leftrightarrow T - \lambda I \text{ is surjective}.\]
\end{lemma}

\begin{proof}
($\Rightarrow$) Suppose $\overline{\lambda} \notin \sigma_{ap}(T^*)$. By definition, there exists $C > 0$ such that
\[\forall x \in K : \|T^*x - \overline{\lambda}x\| \geq C\|x\|.\]
This immediately implies $\Ker(T^* - \overline{\lambda}I) = \{0\}$. \\
Moreover, $\Ran(T^* - \overline{\lambda}I)$ is closed: if $y_n = (T^* - \overline{\lambda}I)(x_n)$ with $y_n \to y$, then
\[\|x_n - x_m\| \leq C^{-1}\|y_n - y_m\|,\]
so $(x_n)$ is Cauchy. Thus $x_n \to x$ for some $x \in K$, and by continuity, $y = (T^* - \overline{\lambda}I)(x)$. \\
By the closed range theorem, $\Ran(T - \lambda I)$ is also closed. Furthermore,
\[\overline{\Ran(T - \lambda I)} = \Ker(T^* - \overline{\lambda}I)^{\perp} = \{0\}^{\perp} = K.\]
Therefore, $T - \lambda I$ is surjective. \\
($\Leftarrow$) Suppose $T - \lambda I$ is surjective. Let $B$ be the Moore-Penrose pseudo-inverse of $T - \lambda I$, defined on $\Ran(T - \lambda I) = K$. Then $B(T - \lambda I) = P_{\Ker(T-\lambda I)^{\perp}}$, the orthogonal projection. \\
For any $x \in K$ with $\|x\| = 1$, write $x = x_1 + x_2$ where $x_1 \in \Ker(T - \lambda I)^{\perp}$ and $x_2 \in \Ker(T - \lambda I)$. Since $B$ exists and is bounded, there exists $\delta > 0$ such that
\[\|x_1\| \leq \delta\|(T - \lambda I)x_1\| = \delta\|(T - \lambda I)x\|.\]
Let $y \in K$ with $\|y\|=1$. By surjectivity of $T-\lambda I$, there exists $x \in K$ such that
\[
(T-\lambda I)x = y.
\]
Decompose $x = x_1 + x_2$ with $x_1 \in \Ker(T-\lambda I)^\perp$ and $x_2 \in \Ker(T-\lambda I)$. Then
\[
(T-\lambda I)x = (T-\lambda I)x_1,
\]
hence
\[
\|x_1\| \le \delta \|y\| = \delta.
\]
Now let $z \in K$. Using the identity
\[
\langle (T-\lambda I)x, z \rangle = \langle x, (T^*-\overline{\lambda}I)z \rangle,
\]
we obtain
\[
|\langle y, z \rangle|
= |\langle (T-\lambda I)x_1, z \rangle|
= |\langle x_1, (T^*-\overline{\lambda}I)z \rangle|
\le \|x_1\| \cdot \|(T^*-\overline{\lambda}I)z\|
\le \delta \|(T^*-\overline{\lambda}I)z\|.
\]
Taking the supremum over all $y$ with $\|y\|=1$, we deduce
\[
\|z\|
= \sup_{\|y\|=1} |\langle y,z\rangle|
\le \delta \|(T^*-\overline{\lambda}I)z\|.
\]
Thus
\[
\|(T^*-\overline{\lambda}I)z\| \ge \delta^{-1}\|z\|
\quad \text{for all } z \in K.
\]
This proves that $T^*-\overline{\lambda}I$ is bounded below, hence
\[
\overline{\lambda} \notin \sigma_{ap}(T^*).
\]
Thus $\overline{\lambda} \notin \sigma_{ap}(T^*)$. 
\end{proof}

\end{document}
